\chapter{Khi Thầy Cô Xin Lỗi}
\markboth{Khi Thầy Cô Xin Lỗi}{Khi Thầy Cô Xin Lỗi}

\section{Người Lớn Sợ Mất Uy Khi Xin Lỗi}
\begin{truyen}{Người Lớn Sợ Mất Uy Khi Xin Lỗi}{Adults Fear Losing Authority When They Apologize}
\chuhoa{H}{ọc sinh:} ``Thầy cô mà xin lỗi học sinh chắc khó lắm ha thầy?''
\textit{(Student: ``It must be hard for teachers to apologize to students, right?'')}

\teacherpair{Khó thật. Vì nhiều người lớn sợ rằng xin lỗi là tự hạ thấp mình. Nhưng thật ra cái làm người dạy mất uy không phải là lời xin lỗi. Cái làm mất uy là sai rõ rồi mà vẫn cố giữ vai đúng.}{``It is hard indeed. Many adults fear that apologizing lowers them. But what truly damages authority is not an apology. It is being obviously wrong and still insisting on playing the role of the one who is right.''}
\end{truyen}

\section{Một Lời Xin Lỗi Dạy Nhiều Hơn Một Bài Giảng}
\begin{truyen}{Một Lời Xin Lỗi Dạy Nhiều Hơn Một Bài Giảng}{One Apology Can Teach More Than a Lecture}
\chuhoa{H}{ọc sinh:} ``Nếu thầy cô xin lỗi tụi em, tụi em có coi thường không?''
\textit{(Student: ``If a teacher apologizes, will students look down on them?'')}

\teacherpair{Có vài đứa non thì có thể tưởng vậy. Nhưng những học sinh tử tế sẽ nhớ rất lâu. Vì lần đầu tiên chúng thấy một người có quyền lực vẫn dám nhận phần sai của mình. Đó là một bài học đạo đức sống còn hơn nhiều bài đạo đức trên bảng.}{``A few immature students may think so. But decent students remember it for a very long time. Because for once they see someone with authority still dare to admit fault. That is a living moral lesson stronger than many moral lessons written on the board.''}
\end{truyen}

\section{Xin Lỗi Không Làm Mất Ranh Giới}
\begin{truyen}{Xin Lỗi Không Làm Mất Ranh Giới}{Apologizing Does Not Destroy Boundaries}
\chuhoa{H}{ọc sinh:} ``Vậy xin lỗi rồi có bị học sinh leo lên đầu không?''
\textit{(Student: ``Then after apologizing, won't students climb onto the teacher's head?'')}

\teacherpair{Không, nếu người xin lỗi vẫn rõ ràng và có bản lĩnh. Xin lỗi là sửa một chỗ sai, không phải tháo luôn kỷ cương. Người lớn yếu là người hoặc không dám nhận lỗi, hoặc sau khi nhận lỗi thì tự tan chảy luôn. Còn người lớn vững là người biết xin lỗi xong vẫn giữ được nguyên tắc.}{``No, not if the apologizing adult remains clear and steady. Apologizing corrects one mistake; it does not erase all discipline. A weak adult either cannot admit fault or collapses completely after doing so. A strong adult apologizes and still keeps the principle intact.''}
\end{truyen}

\begin{baihoc}
Xin lỗi đúng lúc không làm người dạy nhỏ đi.
\textit{(A timely apology does not make a teacher smaller.)}

Nó làm sự tử tế và độ tin cậy của người dạy lớn lên.
\textit{(It makes the teacher's decency and trustworthiness larger.)}
\end{baihoc}

\begin{ghinhoanh}
\textbf{Short English to remember:}

Apology is not weakness.

Authority grows with honesty.
\end{ghinhoanh}

\begin{tuhocnhanh}
1. Em có đang nghĩ người có quyền thì không được sai không? \textit{(Do you assume that people in authority are not allowed to be wrong?)}

2. Nếu em là giáo viên, em có đủ mạnh để xin lỗi học sinh khi cần không? \textit{(If you were a teacher, would you be strong enough to apologize when needed?)}
\end{tuhocnhanh}

\ngancach